% TOP_PROBLEM: {"id": 30007335, "title": "Halting detection under reciprocal-time decaying torus forcing", "category_id": 9, "category": {"id": 9, "name": "pde", "display_name": "Partial Differential Equations", "description": "PDEs and their applications in physics and geometry.", "slug": "pde", "order_index": 9, "created_at": "2026-07-31T15:26:25.671Z"}, "status": "open"}
% FIRST_POSTED: 2026-10-06
% ATTEMPT_STATUS: solved
% OPENAI_CLAIM: {"schema_version":1,"source_repo":"https://github.com/openai/math","source_commit":"adc7f1241b42e322a6451854ab7e4b4c146bf78a","release_date":"2026-10-06","match":"stronger","resolution":"proved","independently_reviewed":false,"second_pass":"agreed","families":[{"family":"376","title":"Universal computation in forced Navier–Stokes flows","subject":"Partial differential equations","manuscripts":[{"title":"Prefix Instructions and Incompressible Flows","pdf_path":"preprints/Prefix-Instructions-and-Incompressible-Flows-September-27-2026/manuscript.pdf","readme_path":"preprints/Prefix-Instructions-and-Incompressible-Flows-September-27-2026/README.md","date":"2026-09-27","theorem_ref":"thm:sa-autonomous (autonomous.tex:10)","statement_sources":[{"path":"preprints/Prefix-Instructions-and-Incompressible-Flows-September-27-2026/build/autonomous.tex","line":10,"label":"thm:sa-autonomous"}]},{"title":"Scalar Potentials and Slow Clocks for Forced Fluid Computation","pdf_path":"preprints/Scalar-Potentials-and-Slow-Clocks-for-Forced-Fluid-Computation-September-27-2026/manuscript.pdf","readme_path":"preprints/Scalar-Potentials-and-Slow-Clocks-for-Forced-Fluid-Computation-September-27-2026/README.md","date":"2026-09-27","theorem_ref":"bcs:prefix-theorem (bcs-prefix.tex:4)","statement_sources":[{"path":"preprints/Scalar-Potentials-and-Slow-Clocks-for-Forced-Fluid-Computation-September-27-2026/build/bcs-prefix.tex","line":4,"label":"bcs:prefix-theorem"}]}],"lean":{"doc_path":"lean/docs/376.md","comparators":[],"declarations":[],"covers_main_theorem":false},"reasoning_trace":null}]}
\documentclass[11pt]{article}
\usepackage[margin=25mm]{geometry}
\usepackage{fontspec}
\setmainfont{Latin Modern Roman}
\usepackage{amsmath,amssymb,mathtools}
\usepackage{unicode-math}
\setmathfont{Latin Modern Math}
\usepackage{enumitem}
\usepackage{xurl,hyperref,bookmark}
\hypersetup{colorlinks=true,urlcolor=blue}
\setcounter{secnumdepth}{0}
\setlength{\parindent}{0pt}
\setlength{\parskip}{5pt}
\setlength{\emergencystretch}{3em}
\begin{document}
\section*{Halting detection under reciprocal-time decaying torus forcing}

\subsection{OpenAI's claim}
\textit{Area-changing prefix instructions}: The manuscript claims complete-domain realization of finite prefix instructions by smooth incompressible motions that retain discarded information; on the flat torus, the resulting mean-zero force detects halting from rest and repeats after loading.

\textit{Scalar potentials and slow-clock computation}: The paper claims compactly supported forced fluid computation with a fixed particle and half-space detector, with time derivatives decaying at the stated inverse-time rates; other realizations include invariant-plane instructions and velocity-field detection under a companion diffusion statement.

\subsection{Mathematical statement of the claimed result}
\paragraph{Mathematical claim 1.} For the incompressible system $\partial_tu+(u\cdot\nabla)u=-\nabla p+\nu\Delta u+f$, $\nabla\cdot u=0$, $u(0)=0$, choose viscosity $\nu>0$, whose value is computable. An effectively computable smooth general body force on $\mathbb T^3$, with spatial mean zero, satisfies \[\|\partial_t^kD_x^\alpha f(t)\|_\infty\le C_{k,\alpha}(1+t)^{-1}\quad(k\ge0).\] The velocity starts from rest and exists smoothly for all times; it is the unique classical periodic solution, its kinetic energy stays bounded and its pressure has spatial mean zero. For the label $y_*=(1/2,1/2,1/4)$ and region $O=\{1/2<x_1<1,\ 0<x_3<1/8\}$, particle entry at some time is equivalent to halting of the machine. The construction records every computation step through the first halt, or forever if it never halts.

\paragraph{Mathematical claim 2.} For the incompressible system $\partial_tu+(u\cdot\nabla)u=-\nabla p+\nu\Delta u+f$, $\nabla\cdot u=0$, $u(0)=0$, choose viscosity $\nu>0$, whose value is computable. The force is smooth and effectively constructed on $\mathbb T^3$, with spatial mean zero. Its corresponding velocity is unique and globally smooth, and both fields satisfy \[\|\partial_t^jD_x^\alpha f(t)\|_\infty+\|\partial_t^jD_x^\alpha u(t)\|_\infty\le C_{j,\alpha}(1+t)^{-1-j}\quad(j\ge0).\] For initial particle position $(1/4,1/2,1/2)$, the event of entering $\{1/2<x_1<1\}$ is equivalent to halting of the supplied machine. A first force choice has zero pressure; its solenoidal projection gives a divergence-free choice with unchanged velocity and altered normalized pressure. Only loading depends on the input for a fixed machine.

\subsection{Outline of the argument}
\textit{Area-changing prefix instructions}: Gapped radix codes turn prefix cylinders into separated closed rectangles. Local curl motions and private heights implement the planar instruction, while reciprocal normal strain compensates its area change. A history invariant and initialization segment identify the simulated state, and horizontal interpolation keeps the detector equivalence valid during the entire instruction period.

\textit{Scalar potentials and slow-clock computation}: Three scalar potentials successively lift a coded rectangle, apply its rule, and lower the image. A logarithmic time change remains onto an unbounded computational-time interval while reducing all time derivatives. Area-changing prefix maps are balanced by a normal component, and the alternative planar processor provides the stated stationary, periodic, and diffusion-based detector variants.

\subsection{Scope relative to this problem}
\textit{Area-changing prefix instructions}: Here only reciprocal-time decay is required for each mixed force derivative, without an extra time power for additional derivatives. The torus construction gives this weaker decay requirement together with bounded kinetic energy and the same fixed particle halting event.

\textit{Scalar potentials and slow-clock computation}: Here only reciprocal-time decay is required for each mixed force derivative, without an extra time power for additional derivatives. The torus construction gives this weaker decay requirement together with bounded kinetic energy and the same fixed particle halting event.

The claimed result implies the stated question and also supplies the additional strength described above.

\subsection{Formal verification}
Family 376 describes Lean coverage for the selected sheet-program, rapidly decaying particle-detector, velocity-detector, and full-box transport statements described separately in its Comparator files.

\textit{Area-changing prefix instructions}: No manuscript-specific Comparator is listed in the pinned inventory.

The documented selection does not cover the entire manuscript assertion relevant here.

\textit{Scalar potentials and slow-clock computation}: No manuscript-specific Comparator is listed in the pinned inventory.

The documented selection does not cover the entire manuscript assertion relevant here.

This site has not run the Lean development. These links describe OpenAI’s selected formal statements and should be read with their stated scope.

\subsection{Status caveat}
OpenAI reports that an internal model produced the manuscripts. Their arguments have not been independently verified as proofs by this site. Any portion outside the documented formal scope may contain mistakes, and the catalogue label reports the announced mathematical claim.

\subsection{Sources}
\begin{itemize}
\item OpenAI, \emph{Prefix Instructions and Incompressible Flows}, 2026-09-27; citation key \texttt{OAI:Prefix-Instructions-and-Incompressible-Flows-September-27-2026}. \href{https://github.com/openai/math/blob/adc7f1241b42e322a6451854ab7e4b4c146bf78a/preprints/Prefix-Instructions-and-Incompressible-Flows-September-27-2026/manuscript.pdf}{PDF (GitHub)}; \href{https://github.com/openai/math/raw/adc7f1241b42e322a6451854ab7e4b4c146bf78a/preprints/Prefix-Instructions-and-Incompressible-Flows-September-27-2026/manuscript.pdf}{PDF (direct)}; \href{https://github.com/openai/math/blob/adc7f1241b42e322a6451854ab7e4b4c146bf78a/preprints/Prefix-Instructions-and-Incompressible-Flows-September-27-2026/README.md}{Manuscript page and citation}; \href{https://github.com/openai/math/blob/main/preprints/Prefix-Instructions-and-Incompressible-Flows-September-27-2026/manuscript.pdf}{Latest version}; Statement source: \href{https://github.com/openai/math/blob/adc7f1241b42e322a6451854ab7e4b4c146bf78a/preprints/Prefix-Instructions-and-Incompressible-Flows-September-27-2026/build/autonomous.tex\#L10}{thm:sa-autonomous}.
\item OpenAI, \emph{Scalar Potentials and Slow Clocks for Forced Fluid Computation}, 2026-09-27; citation key \texttt{OAI:Scalar-Potentials-and-Slow-Clocks-for-Forced-Fluid-Computation-September-27-2026}. \href{https://github.com/openai/math/blob/adc7f1241b42e322a6451854ab7e4b4c146bf78a/preprints/Scalar-Potentials-and-Slow-Clocks-for-Forced-Fluid-Computation-September-27-2026/manuscript.pdf}{PDF (GitHub)}; \href{https://github.com/openai/math/raw/adc7f1241b42e322a6451854ab7e4b4c146bf78a/preprints/Scalar-Potentials-and-Slow-Clocks-for-Forced-Fluid-Computation-September-27-2026/manuscript.pdf}{PDF (direct)}; \href{https://github.com/openai/math/blob/adc7f1241b42e322a6451854ab7e4b4c146bf78a/preprints/Scalar-Potentials-and-Slow-Clocks-for-Forced-Fluid-Computation-September-27-2026/README.md}{Manuscript page and citation}; \href{https://github.com/openai/math/blob/main/preprints/Scalar-Potentials-and-Slow-Clocks-for-Forced-Fluid-Computation-September-27-2026/manuscript.pdf}{Latest version}; Statement source: \href{https://github.com/openai/math/blob/adc7f1241b42e322a6451854ab7e4b4c146bf78a/preprints/Scalar-Potentials-and-Slow-Clocks-for-Forced-Fluid-Computation-September-27-2026/build/bcs-prefix.tex\#L4}{bcs:prefix-theorem}.
\item Family 376: \href{https://github.com/openai/math/blob/adc7f1241b42e322a6451854ab7e4b4c146bf78a/lean/docs/376.md}{Lean scope}.
\item \href{https://github.com/openai/math/blob/adc7f1241b42e322a6451854ab7e4b4c146bf78a/CONTENTS.md}{OpenAI catalogue entry}.
\end{itemize}
\end{document}
